\section{Results}\label{sec:results}

\subsection{Consensus surprises forecast the drift}

We first evaluate the surprise approach on its own, walk-forward over 2019--2026 at our base cost
of two ticks plus \$4.50. Trading the sign of the CPI headline surprise earns a net Sharpe ratio
of 1.03 in \ES\ and 1.00 in \NQ, with 9.8 and 13.6 bps per trade over about 86 and 79 releases. The
CPI model, which trades only when the predicted move exceeds the cost, earns 0.99 in \ES\ and 0.79
in \NQ, with hit rates of 64\% and 58\%. The same rule on the E-mini Dow, a market added after the
\ES\ and \NQ\ results were known, earns 0.99 (77 trades, hit rate 57\%). The 90\% block-bootstrap confidence interval of the \ES\
Sharpe ratio is 0.25 to 1.72. In the format common among practitioners, the \ES\ CPI strategy gains
6.6\% of notional while being long in exactly the same 30-minute windows loses 6.0\%; in \NQ\ the
figures are $+7.5\%$ and $-5.4\%$. The profit therefore comes from the direction of the surprise,
not from being in the market at release times. In \ZN\ the naive rule earns 2.7 bps per trade
before costs (gross Sharpe 0.53) and $-0.3$ bps after them; the CPI model correctly declines to
trade.

Trading the ten screened families rather than CPI alone raises the number of trades tenfold, to 98
a year in \ES\ and 97 in \NQ, with net Sharpe ratios of 0.65 and 0.71, 1.7 and 2.4 bps per trade, and
random-sign $p$-values below 0.001 and of 0.005. CPI remains the largest contributor in both markets
($+4.3\%$ and $+6.0\%$ of notional); behind it the ranking of families differs between \ES\ and \NQ,
so the broad book should be read as diversification across many small edges. On \YM\ the broad book
earns 0.68 (about 93 trades a year, random-sign $p=0.006$). Combining \ES\ and \NQ\ gives about 195
trades a year and a Sharpe ratio of 0.72 (random-sign $p<0.001$); adding \YM\ gives about 290 trades
and 0.78, with a 90\% confidence interval of 0.10 to 1.50.

\subsection{Transaction costs and turnover}

Figure~\ref{fig:costs} shows how the CPI strategy's Sharpe ratio changes as the round-trip
slippage rises from zero to eight ticks. The line falls gently in \ES\ (1.15 to 0.62) and is almost
flat in \NQ\ (1.02 to 0.81), because one \NQ\ tick is only 0.1--0.3 bps of the contract value. \YM,
whose tick is about 0.25--1 bp, sits between them (1.06 to 0.81). In \ZN\ the naive rule crosses
zero at about two ticks. For the broad book, the \ES\ Sharpe ratio falls from 1.09 gross to 0.20 at
four ticks, while \NQ\ stays at 0.60 and \YM\ falls from 0.91 to 0.45. Cost robustness is therefore a
property of the market's tick size relative to its post-release moves, not only of the strategy.

\begin{figure}[htbp]
  \centering
  \includegraphics[width=\linewidth]{fig20_cost_sensitivity_4markets}
  \caption{Net Sharpe ratio of the CPI strategies as round-trip slippage rises from 0 to 8 ticks
  (plus a \$4.50 commission), \ES, \NQ, \YM\ and \ZN, 2019--2026.}
  \label{fig:costs}
\end{figure}

\subsection{Four approaches on equal terms}

Table~\ref{tab:strategies} compares the four approaches on the held-out years 2021--2026 with
the same bars, entry, costs and metrics. At a quarter tick per side, the surprise book has the
highest Sharpe ratio in \ES\ (1.09), \NQ\ (0.75) and \YM\ (1.06), followed in \ES\ by the classifier
(0.95), the DTW ladder (0.64) and the event games (0.42). At our base cost the ranking changes: the
classifier (0.81) and the surprise book (0.74) remain clearly positive, the DTW ladder falls to
0.10, and the event games, which trade about 370 times a year, turn negative. In \ZN\ every
approach loses money after costs, although the DTW ladder earns a gross Sharpe ratio of 0.34. On
\YM\ the DTW ladder, run with the same code and setting, loses even before costs ($-0.32$ gross,
$-0.39$ at a quarter tick), while the surprise book earns 1.06: the third equity index confirms the
surprise and not the price-path analogues. The
surprise book is positive in both the validation and the test years on every equity market (\ES\
1.04 and 1.09, \NQ\ 0.86 and 0.75, \YM\ 1.27 and 1.06).

\begin{table}[htbp]
  \centering\small
  \caption{Sharpe ratio of each strategy on each market, test years 2021--2026. ``--'' = the
  strategy is not available on that market.}
  \label{tab:strategies}
  \begin{threeparttable}
  \setlength{\tabcolsep}{4pt}
  \begin{tabular}{lccccccccc}
    \toprule
    & \multicolumn{4}{c}{$\tfrac14$ tick per side} & \multicolumn{4}{c}{2 ticks + \$4.50} & Valid. \\
    \cmidrule(lr){2-5}\cmidrule(lr){6-9}\cmidrule(lr){10-10}
    Strategy & \ES & \NQ & \ZN & \YM & \ES & \NQ & \ZN & \YM & \ES \\
    \midrule
    Macro-surprise book         & 1.09 & 0.75 & $-0.16$ & 1.06 & 0.74 & 0.64 & $-1.51$ & 0.82 & 1.04 \\
    CPI surprise model          & 0.51 & 0.46 & 0.17    & 0.39 & 0.61 & 0.45 & $-0.43$ & 0.29 & 1.27 \\
    DTW ladder book             & 0.64 & 0.11 & $-0.97$ & $-0.39$ & 0.10 & $-0.04$ & $-5.03$ & $-0.74$ & 0.59 \\
    DTW event-game book         & 0.42 & -- & -- & -- & $-0.29$ & -- & -- & -- & 0.26 \\
    Machine-learning classifier & 0.95 & -- & -- & -- & 0.81 & -- & -- & -- & 0.38 \\
    \midrule
    Buy and hold (no costs)     & 0.69 & 0.58 & $-0.69$ & 0.67 & & & & & 0.80 \\
    \bottomrule
  \end{tabular}
  \begin{tablenotes}\footnotesize
    \item Notes: Daily Sharpe ratios on the session grid. The DTW ladder uses the authors' own feature
    panel and rule; the validation column reports 2018--2020 at $\tfrac14$ tick per side.
  \end{tablenotes}
  \end{threeparttable}
\end{table}


\subsection{Combining surprises and price-path analogues}

If the price path before the entry carried information beyond the surprise, conditioning the
surprise on the analogues should improve it. It does not. Across all releases from 2019 to 2026, the
rank correlation of the DTW direction forecast with the next 30-minute return is 0.010 in \ES,
0.026 in \NQ, 0.030 in \YM\ and 0.039 in \ZN, its hit rate is 47--51\%, and its correlation with the
surprise is 0.02--0.13. In the held-out comparison (Figure~\ref{fig:surprise_dtw}), keeping a surprise trade only
when the analogues agree raises the \ES\ Sharpe ratio in the validation years (1.46 against 1.04)
but lowers it in the test years (0.37 against 1.09); \NQ\ shows the same reversal (1.21 against 0.86,
then 0.48 against 0.75). The reason is that the direction of the analogues itself changed sign: the
DTW-only book earned 0.57 (\ES) and 0.78 (\NQ) in 2018--2020 and $-0.26$ and $-0.28$ in 2021--2026.
\YM\ does not show the reversal, but it does not show a gain either: its DTW-only book earns 0.40
and 0.44 in the two periods, and the agreement filter lowers the surprise book in the validation
years (0.63 against 1.27) and leaves it unchanged in the test years (1.04 against 1.06). In the
walk-forward evaluation on 2019--2026 a blend of the two signals raised the \YM\ broad book from
0.68 to 1.06, but the bootstrap probability that it is not better is 6\%, and the same blend does
not help in \ES\ or \NQ. A filter that helps in only one of two periods, or on one of three equity
indices, is not a reliable source of information. In
contrast, the dispersion of the analogues predicts the size of the next move with a rank
correlation of 0.10--0.30 in every market and year, and sizing surprise trades inversely to it
gives Sharpe ratios of 0.63 and 0.72 in \ES\ and 0.64 and 0.97 in \NQ\ in the training and validation
years.

\begin{figure}[htbp]
  \centering
  \includegraphics[width=\linewidth]{surprise_dtw_book_cumulative}
  \caption{Cumulative net P\&L at a quarter tick per side, 2021--2026, of the surprise books
  (top-10 families and CPI model), the DTW-only book and the surprise books filtered by DTW
  agreement, in \ES, \NQ, \ZN, \YM\ and equal-risk \ES\ + \NQ\ and \ES\ + \NQ\ + \YM\ books.}
  \label{fig:surprise_dtw}
\end{figure}

\subsection{The entry assumption}

Table~\ref{tab:equalcosts} re-prices each approach's own \ES\ book on 2019--2026 at three cost
levels. The classifier's reported Sharpe ratio of 2.14 (total return 38.6\% gross) falls to 0.92
(7.3\%) when the same positions are entered one minute later: about 80\% of the reported profit is
the first-minute jump. An independent rebuild of the classifier's pipeline finds a loss of about
96\% of the profit per trade from the same one-minute delay. The first minute is where the market
prices the news, and it cannot be traded.

\begin{table}[htbp]
  \centering\small
  \caption{Each workstream's own book re-priced at three cost levels, \ES, 2019--2026
  (Sharpe / total return).}
  \label{tab:equalcosts}
  \begin{tabular}{lccc}
    \toprule
    Book & Gross & $\tfrac14$ tick & 2 ticks + \$4.50 \\
    \midrule
    Macro-surprise, all releases & 1.23 / 23.3\% & 1.12 / 21.1\% & 0.69 / 13.0\% \\
    DTW ladder (2018 \ES\ rebuild)$^{a}$ & 0.27 / 8.8\%  & 0.07 / 2.4\%  & $-0.66$ / $-21.5\%$ \\
    DTW event games              & 0.45 / 18.6\% & 0.25 / 10.3\% & $-0.50$ / $-20.6\%$ \\
    Classifier, executable       & 0.92 / 7.3\%  & 0.88 / 6.9\%  & 0.72 / 5.6\% \\
    Classifier, as reported      & 2.14 / 38.6\% & 2.13 / 38.3\% & 2.07 / 37.0\% \\
    \bottomrule
  \end{tabular}
  \par\smallskip\footnotesize $^{a}$ Rebuilt from the published 2018 specification; the current
  book is in Table~\ref{tab:strategies}.
\end{table}


\subsection{Diversification across approaches}

Daily returns of the four approaches are nearly uncorrelated. In the test years the correlations
across approaches range from $-0.04$ to 0.31, the highest being between the surprise book and the
classifier, which also uses surprise features; the surprise and DTW ladder books correlate at 0.13
on \ES. The high correlations are the same surprise strategy on different equity indices (\ES\ and
\NQ\ 0.68, \ES\ and \YM\ 0.85), because the same releases move all of them. We therefore combine
the approaches with equal-risk weights, proportional to the inverse of each book's daily volatility
in 2018--2020, and score them on 2021--2026 (Table~\ref{tab:portfolios} and Figure~\ref{fig:team}).
At a quarter tick per side, the \ES\ portfolio of the four approaches earns a Sharpe ratio of
\textbf{1.39}, the equity portfolio on \ES, \NQ\ and \YM\ 1.35 and the \ES\ + \NQ\ portfolio 1.31, above
the best single approach (1.09) with less than half its volatility (1.2--1.3\% against 2.7\%) and a
maximum drawdown of about $-1\%$. The surprise book and the DTW ladder alone earn 1.09 and 0.64 on
\ES; together they earn 1.16, because they lose in different years: 2022, the surprise book's best
year, is the only year in which the DTW ladder loses. At our base cost the equity portfolio on
\ES, \NQ\ and \YM\ holds up best (0.85), ahead of the surprise book on the same three markets
(0.82) and of the \ES\ portfolio (0.76), whose high-turnover pattern-based books pay the cost on 300
or more trades a year. Adding \ZN\ lowers the combined Sharpe ratio to 0.83 at a quarter tick and
turns it negative at our base cost.

\begin{table}[htbp]
  \centering\small
  \caption{Equal-risk team portfolios, test years 2021--2026 (weights from 2018--2020).
  Return, volatility and drawdown at $\tfrac14$ tick per side.}
  \label{tab:portfolios}
  \setlength{\tabcolsep}{4pt}
  \begin{tabular}{lcccccccc}
    \toprule
    & \multicolumn{4}{c}{Sharpe} & & & & \\
    \cmidrule(lr){2-5}
    Portfolio & Gross & $\tfrac14$ tick & $\tfrac12$ tick & 2t+\$4.50 & Return & Vol & Max DD & Trades/yr \\
    \midrule
    Team, \ES\ (four strategies)       & 1.55 & 1.39 & 1.22 & 0.76    & 8.9\%  & 1.2\% & $-1.0\%$ & 739 \\
    Team, equities (\ES+\NQ+\YM)       & 1.48 & 1.35 & 1.22 & 0.85    & 9.6\%  & 1.3\% & $-1.1\%$ & 1{,}131 \\
    Team, \ES+\NQ                      & 1.45 & 1.31 & 1.17 & 0.78    & 9.0\%  & 1.3\% & $-1.1\%$ & 1{,}043 \\
    Team, all markets (with \ZN)       & 1.35 & 0.83 & 0.30 & $-1.09$ & 4.3\%  & 1.0\% & $-1.1\%$ & 1{,}562 \\
    Surprise + DTW ladder (\ES)        & 1.31 & 1.16 & 1.01 & 0.59    & 13.8\% & 2.2\% & $-2.0\%$ & 315 \\
    Surprise + DTW event games (\ES)   & 1.32 & 1.14 & 0.96 & 0.46    & 14.8\% & 2.4\% & $-2.2\%$ & 465 \\
    \midrule
    Macro-surprise book, \ES+\NQ+\YM   & 1.13 & 1.07 & 1.01 & 0.82    & 14.6\% & 2.6\% & $-2.4\%$ & 280 \\
    Macro-surprise book, \ES\ alone    & 1.18 & 1.09 & 0.99 & 0.74    & 15.9\% & 2.7\% & $-2.6\%$ & 96 \\
    DTW ladder, \ES+\NQ\ (own weights) & 0.66 & 0.54 & 0.42 & 0.07    & 7.4\%  & 2.6\% & $-5.3\%$ & 427 \\
    \bottomrule
  \end{tabular}
\end{table}


\begin{figure}[htbp]
  \centering
  \includegraphics[width=\linewidth]{team_report_figure}
  \caption{Cumulative net P\&L at a quarter tick per side, 2021--2026: each approach in \ES\
  (left), in \NQ, \ZN\ and \YM\ (center), and the team portfolios with the DTW ladder's own
  \ES\ + \NQ\ book (right).}
  \label{fig:team}
\end{figure}
